% An inertial baseline over the whole corpus, and an accounting of what that baseline
% takes from the camera reference. Every number here is produced by the scripts named
% below; nothing is hand-entered. Regenerate with:
%     python3 scripts/imu/pipeline.py      --n 84     % Table \ref{tab:imu_baseline}, top
%     python3 scripts/imu/bar_path.py      --n 84     % Table \ref{tab:imu_baseline}, foot
%     python3 scripts/imu/independence.py  --n 84     % Table \ref{tab:imu_ablation}
%     python3 scripts/imu/compare_attitude.py --n 20  % attitude-filter paragraph
% The per-session figures are the audit images released with the corpus,
% <session>/audit_imu.png, produced by scripts/imu/audit_imu.py --n 84.

\section{An Inertial Baseline, and What It Owes the Reference}\label{sec:imu_baseline}

The corpus exists so that inertial estimation can be trained and judged against something
that is not itself inertial. This section establishes a baseline on it, and then does
something a velocity-estimation result is not usually asked to do: it accounts for what
the estimate takes from the reference, and measures how much accuracy survives when each
of those things is taken away. The motivation is that a figure obtained inside a
repetition whose start and end were supplied by a camera is not a figure a bar-mounted
device can reproduce, and the difference between the two has, to our knowledge, not been
quantified on a corpus of this size.

\subsection{What the Estimate Takes From the Reference, and What It Does Not}
\label{sec:imu_given}

Three things enter the estimator from the camera side, and they are stated here rather
than left implicit.

\paragraph{(i) Repetition boundaries.} Each repetition's first and last frame come from
the released annotation. Inside the estimator they do three jobs: they choose the
integration window, they impose $v(0) = v(T) = 0$ on the vertical velocity, and they
impose $p(0) = p(T) = 0$ on the vertical displacement. The concentric interval, also from
the annotation, chooses the window over which peak and mean velocity are read.

The two conditions are earned by the geometry of a repetition rather than assumed. A
boundary is the far end of a round trip, so the vertical velocity there really is near
zero --- the camera measures \SI{33}{\milli\metre\per\second} at a boundary against a
\SI{1050}{\milli\metre\per\second} peak --- and the bar returns to where it set off, which
the camera puts within \SI{15}{\milli\metre} on a \SI{490}{\milli\metre} travel. The
velocity residual is removed as a ramp, the shape a constant accelerometer bias or a
constant gravity leak would produce; the displacement residual is then removed as the
parabola that carries the displacement back without moving either endpoint.

\paragraph{(ii) A lever arm.} The inertial sensor sits on the left collar and the optical
marker elsewhere on the bar, so the two points differ by $\bm{\omega} \times \bm{r}$. A
single $|\bm{r}| = \SI{12.3}{\centi\metre}$ was fitted once against the camera and used
for the whole corpus. Section~\ref{sec:imu_ablation} shows this is a convenience rather
than a dependency.

\paragraph{(iii) One heading per session, for the path only.} A six-axis inertial sensor
cannot recover heading: rotating the world about the vertical changes nothing it measures.
One yaw angle per session was therefore fitted against the camera, and the horizontal and
three-dimensional path figures are labelled accordingly. The height, the path length and
the departure from vertical need no heading and are reported separately.

\paragraph{What is not taken.} There is no zero-velocity update anywhere in the pipeline,
and none is possible: at a repetition boundary the bar is still rotating at
\SI{14.2}{\degree\per\second} (median over the corpus), because a person holding a loaded
barbell never stops moving it. Anything anchored on detected stillness is ruled out, which
is precisely why the round-trip conditions above are used instead. Gyroscope bias and
accelerometer scale come from windows \emph{found} to be still in the sensor's own angular
rate --- never from a position in the session, since the interval before the first
repetition is not still, the bar being handled. The time base comes from the hardware pulse
train of Section~\ref{sec:sync_measured}, not from the camera's position signal. The
turnaround instant, which the annotation also records, is not used at all.

\subsection{Baseline Accuracy Given Reference Boundaries}\label{sec:imu_accuracy}

Table~\ref{tab:imu_baseline} reports the baseline over all \num{84} sessions and
\num{1400} repetitions. Orientation is from VQF~\cite{laidig2023}; the vertical
specific force is taken in the world frame with gravity removed, band-limited at
\SI{10}{\hertz}, and integrated one repetition at a time under the conditions of
Section~\ref{sec:imu_given}.

\begin{table}[t]
  \centering
  \caption{Inertial estimation against the camera reference, all \num{84} sessions and
           \num{1400} repetitions, with the repetition boundaries taken from the released
           annotation. Velocity and range of motion are compared repetition by
           repetition; the path figures are the median over repetitions of the RMS error
           within a repetition.}
  \label{tab:imu_baseline}
  \begin{tabular}{lrrrr}
    \toprule
    & RMSE & bias & 95\% LoA & $R^2$ \\
    \midrule
    \multicolumn{5}{l}{\emph{Per repetition, against the camera --- attitude from VQF}} \\
    Peak concentric velocity (\si{\milli\metre\per\second})  & 50.5 & $-10.9$ & $\pm 96.7$  & 0.96 \\
    Mean concentric velocity (\si{\milli\metre\per\second})  & 36.0 & $-6.6$  & $\pm 69.3$  & 0.95 \\
    Range of motion (\si{\milli\metre})                      & 55.7 & $-17.5$ & $\pm 103.7$ & 0.79 \\
    \midrule
    \multicolumn{5}{l}{\emph{The same, with attitude from the error-state filter of Sec.~\ref{sec:imu_accuracy}}} \\
    Peak concentric velocity (\si{\milli\metre\per\second})  & 49.9 & $-11.7$ & $\pm 95.0$  & 0.96 \\
    Mean concentric velocity (\si{\milli\metre\per\second})  & 35.1 & $-7.5$  & $\pm 67.1$  & 0.95 \\
    Range of motion (\si{\milli\metre})                      & 54.9 & $-18.4$ & $\pm 101.5$ & 0.79 \\
    \midrule
    & median & 90th & & \\
    \midrule
    \multicolumn{5}{l}{\emph{Path, no heading required}} \\
    Height over the repetition (\si{\milli\metre})           & 20.1 & 57.5 & & \\
    Length of the path (\si{\milli\metre}, RMSE)             & 140.4 & & \multicolumn{2}{l}{\footnotesize camera median 1063} \\
    Furthest departure from vertical (\si{\milli\metre}, RMSE) & 67.2 & & \multicolumn{2}{l}{\footnotesize camera median 114} \\
    \midrule
    \multicolumn{5}{l}{\emph{Path, after borrowing one heading per session}} \\
    Horizontal (\si{\milli\metre})                           & 36.5 & 77.9 & & \\
    Three-dimensional (\si{\milli\metre})                    & 45.9 & 97.9 & & \\
    \bottomrule
  \end{tabular}
\end{table}

The height is the quantity the round-trip conditions act on directly, and it is the one
the inertial sensor recovers best: \SI{20.1}{\milli\metre} of RMS error over a repetition,
against a median travel above \SI{490}{\milli\metre}, or about \SI{4}{\percent}. The
horizontal is worse and for an identifiable reason. One degree of tilt error puts
\SI{0.171}{\metre\per\second\squared} into the horizontal, which over a two-second
repetition integrates to \SI{342}{\milli\metre}; the horizontal path error on this corpus
tracks repetition duration at $+0.44$ and rotation rate at $+0.35$, which is the signature
of exactly that. The vertical does not suffer it because gravity defines the vertical and
the boundary conditions absorb what is left.

\paragraph{The attitude filter is not the limiting term for velocity.} An error-state
Kalman filter on $[\,$tilt error, gyroscope bias$\,]$ and its iterated form were
implemented and compared against VQF with everything else held identical --- same
calibration, same lever arm, same band-limit, same conditions, same repetitions. On
\num{294} repetitions the three agree on velocity to within
\SI{0.2}{\milli\metre\per\second} on peak. They separate only on the path, where VQF leads
by about \SI{4}{\milli\metre} of horizontal error (\num{32.5} against
\SI{37.0}{\milli\metre}). Iterating the update changes nothing measurable: at
\SI{1}{\kilo\hertz} the per-sample correction is far too small for the nonlinearity of the
gravity direction to bite, so the iterated filter reproduces the single-pass filter at
every iteration count tried. This is a negative result worth recording, because it locates
the remaining error in the sensor and the geometry rather than in the choice of estimator.

\subsection{How Much the Reference Boundaries Are Worth}\label{sec:imu_ablation}

Table~\ref{tab:imu_ablation} takes the camera away a piece at a time. Displaced boundaries
are \emph{scored}, never discarded: dropping the repetitions a mis-placed boundary ruined
is how a fragile method comes to look robust. In the boundary-free rows the round-trip
conditions are replaced by a high-pass on the velocity, drift being the lowest-frequency
thing present; the high-pass is applied to the velocity and not to the acceleration, since
the drift is created by the integration and is only there to be removed afterwards.

\begin{table}[t]
  \centering
  \caption{RMS error of concentric velocity as the reference is withdrawn. All \num{84}
           sessions and \num{1400} repetitions in every row. One camera frame is
           \SI{11.1}{\milli\second}.}
  \label{tab:imu_ablation}
  \begin{tabular}{lrr}
    \toprule
    What the camera supplies & Peak (\si{\milli\metre\per\second}) & Mean (\si{\milli\metre\per\second}) \\
    \midrule
    Boundaries, exact                        & 50.5   & 36.0 \\
    \midrule
    \multicolumn{3}{l}{\emph{Both boundaries displaced the same way (a systematic offset)}} \\
    \quad by $\pm 2$ frames                 & 52.6   & 37.9 \\
    \quad by $\pm 5$ frames                 & 61.1   & 46.7 \\
    \quad by $\pm 10$ frames                & 88.4   & 80.1 \\
    \quad by $\pm 20$ frames                & 255.9  & 207.5 \\
    \multicolumn{3}{l}{\emph{Displaced oppositely (the repetition stretched or squeezed)}} \\
    \quad by $\pm 2$ frames                 & 60.5   & 48.0 \\
    \quad by $\pm 5$ frames                 & 90.3   & 78.4 \\
    \quad by $\pm 10$ frames                & 140.4  & 145.1 \\
    \quad by $\pm 20$ frames                & 232.2  & 274.0 \\
    \multicolumn{3}{l}{\emph{Each boundary drawn independently}} \\
    \quad by $\pm 2$ frames                 & 58.5   & 45.2 \\
    \quad by $\pm 10$ frames                & 114.7  & 113.0 \\
    \midrule
    The start only, no round trip            & 92.4   & 91.3 \\
    \midrule
    Nothing; high-pass at \SI{0.1}{\hertz}   & 56.9   & 51.0 \\
    Nothing; no high-pass                    & 1406.4 & 1388.4 \\
    \midrule
    \multicolumn{3}{l}{\emph{Boundaries exact, lever arm varied}} \\
    No lever arm                             & 70.1   & 42.9 \\
    Lever arm at half                        & 53.2   & 36.4 \\
    Lever arm at double                      & 87.4   & 52.0 \\
    \bottomrule
  \end{tabular}
\end{table}

The displaced-boundary rows are single random draws; re-drawing moves them by one to four
\si{\milli\metre\per\second}, so differences of that size between adjacent rows are draw
noise rather than signal.

Four things follow, and the first is not what we expected.


\paragraph{The boundaries are not what buys the accuracy; precise boundaries are.}
Removing them entirely and high-passing instead costs
\SI{6}{\milli\metre\per\second} on peak velocity (\num{50.5} to
\num{56.9}). Keeping them but placing them ten frames off costs between \num{38} and
\SI{90}{\milli\metre\per\second} more, depending on how that error is distributed
between the two ends. A boundary condition imposed at the wrong instant is not a weak
constraint but a wrong one: it injects into the velocity a ramp that was never there.

\paragraph{Consistency matters more than accuracy, by a factor of two and a half.}
Displacing both boundaries the same way is partly cancelled by the velocity-closure ramp,
which removes a common error exactly; the residual costs about
\SI{3.8}{\milli\metre\per\second} per camera frame. Displacing them oppositely stretches
the repetition, which the ramp cannot cancel, and costs about
\SI{9.0}{\milli\metre\per\second} per frame. Measured against the boundary-free figure
this gives two tolerances rather than one: a purely systematic offset is affordable out to
three or four frames (\SI{35}{} to \SI{45}{\milli\second}), while inconsistency between
the two ends of the same repetition must stay inside roughly one and a half frames
(\SI{17}{\milli\second}). A detector should therefore be built to lock onto a repeatable
feature of the signal rather than to estimate the true turnaround: low jitter with a known
systematic offset is worth more here than an unbiased estimate with scatter. We are not
aware of this split being reported elsewhere, and it changes the design goal from accuracy
to repeatability.

\paragraph{The lever arm has to be present and does not have to be right.} Omitting it
costs \SI{20}{\milli\metre\per\second} on peak velocity, so the rigid-body term between
sensor and marker is real and cannot be neglected. But halving or doubling it costs
between \num{3} and \SI{37}{\milli\metre\per\second}, and halving it is within
\SI{3}{\milli\metre\per\second} of the fitted value. The
\SI{12.3}{\centi\metre} obtained against the camera can therefore be replaced by a tape
measure against the mounting bracket without material loss. It is not a dependency on the
reference; it was a convenience.

\paragraph{Withdrawing the reference from the estimator costs little.} The
boundary-free row, \SI{56.9}{\milli\metre\per\second} on peak and
\num{51.0} on mean, is obtained with no camera-derived quantity inside the estimator at
all. The mean degrades proportionally more than the peak, which is consistent with the
round-trip displacement condition being the term that pins the mean; the peak is set by the
shape of the concentric, which the accelerometer measures directly.

\subsection{What the Boundary-Free Figure Does and Does Not Claim}
\label{sec:imu_limits}

Two limits are stated so the figure is not read for more than it supports.

First, the high-pass is a zero-phase filter run forwards and backwards. That is legitimate
for an offline reference computed on a stored session, which is what this corpus is for,
but a causal device cannot have it, and a causal high-pass would do worse. The
\SI{56.9}{\milli\metre\per\second} is therefore an offline boundary-free figure, not a
claim about what a real-time implementation achieves.

Second, the interval over which the score is computed is still the camera's, in every row
of Table~\ref{tab:imu_ablation} including the boundary-free ones. Something has to say
where the concentric was for the comparison to mean anything. This is the reference used as
a ruler rather than as an input, and the distinction is the point of the table: the
estimator can be made independent of the camera, while the \emph{evaluation} cannot be, and
a device that reports per-repetition velocity would still need to detect repetitions on its
own. That is a counting problem rather than an integration problem, and the two should not
be reported as one number, which is what Table~\ref{tab:imu_ablation} exists to prevent.
